%======================================================================
\section{Vector spaces and linear maps}\label{sec:vector-spaces}
%======================================================================

\begin{objectives}
  You will decide whether a set is a vector space, find bases and dimensions, recognise linear maps and compute their kernel and range, write a linear map as a matrix and change its basis, and use the Pauli matrices, the first quantum operators we meet.
\end{objectives}

The Introduction stressed that Schr\"odinger's equation is linear: its solutions can be added and multiplied by complex numbers, and that is why amplitudes add. The natural setting for such objects is a vector space, and the maps that respect addition and scaling are the linear maps. A great deal of quantum mechanics is just the study of these two notions, so we spend some time on them.

\subsection{Vector spaces}

Throughout, the scalars $F$ will be either $\R$ or $\C$; in quantum mechanics it is nearly always $\C$.

\begin{definition}[Vector space]\label{def:vs}
  A \term{vector space} over $F$ is a set $V$ with an addition $V\times V\to V$ and a scalar multiplication $F\times V\to V$ such that $(V,+)$ is a commutative group, with neutral element $0_V$, and for all $\lambda,\mu\in F$ and $v,w\in V$:
  \[
    \lambda(\mu v)=(\lambda\mu)v,\quad (\lambda+\mu)v=\lambda v+\mu v,\quad \lambda(v+w)=\lambda v+\lambda w,\quad 1\,v=v .
  \]
\end{definition}

Elements of $V$ are called \term{vectors}. In words: you may add vectors, scale them, and the usual rules of arithmetic hold. From these axioms one deduces that $0\,v=0_V$ for every $v$.

The most important example is $F^n$, the columns of $n$ scalars. But the definition covers much more. The $m\times n$ matrices form a vector space, and so do the polynomials of degree at most $n$, written $P_n$. In infinite dimension we have the continuous functions on $[0,1]$, and, because the equation is linear, the solutions of the Schr\"odinger equation. Sequences $(z_1,z_2,\dots)$ of complex numbers with termwise operations form a vector space as well.

To see how the definition is used, and how it fails, consider two subsets of $\R^2$.

\begin{example}\label{ex:vs-nonexample}
  Are $A=\{(x,y)\in\R^2: x+y=0\}$ and $B=\{(x,y)\in\R^2:x+y=1\}$ vector spaces with the usual operations?
\end{example}
\begin{solution}
  \begin{steps}
    \item \emph{$A$ contains $0$.} $(0,0)\in A$.
    \item \emph{$A$ is closed under sums.} If $x+y=0$ and $x'+y'=0$ then $(x+x')+(y+y')=0$.
    \item \emph{$A$ is closed under scalars.} $\lambda x+\lambda y=\lambda(x+y)=0$. The axioms are inherited from $\R^2$, so $A$ is a vector space (a line through the origin).
    \item \emph{$B$ has no zero.} $(0,0)\notin B$, so $B$ cannot be a vector space.
    \item \emph{$B$ fails closure.} $(1,0)$ and $(0,1)$ are in $B$, but their sum $(1,1)$ is not.
  \end{steps}
\end{solution}

The test used for $A$ is the one to remember: to check that a subset of a known vector space is itself a vector space, it is enough to check that it contains $0$ and is closed under sums and scalar multiples. Such a subset is called a \term{subspace}.

\begin{example}[Is the set of $2\times2$ trace-zero matrices a subspace?]\label{ex:trace-zero}
  The \term{trace} of a square matrix is the sum of its diagonal entries, $\Tr M=\sum_iM_{ii}$. Decide whether $W=\{M\in M_2(\C):\Tr M=0\}$ is a subspace of $M_2(\C)$, and find its dimension.
\end{example}
\begin{solution}
  \begin{steps}
    \item \emph{Subspace test.} $\Tr0=0$; and $\Tr(M+N)=\Tr M+\Tr N=0$, $\Tr(\lambda M)=\lambda\Tr M=0$ for $M,N\in W$. So $W$ is a subspace.
    \item \emph{Describe its elements.} $M=\begin{pmatrix}a&b\\c&-a\end{pmatrix}=a\begin{pmatrix}1&0\\0&-1\end{pmatrix}+b\begin{pmatrix}0&1\\0&0\end{pmatrix}+c\begin{pmatrix}0&0\\1&0\end{pmatrix}$.
    \item \emph{Dimension.} The three matrices are independent and generate $W$, so $\dim W=3$ (compare with $\dim M_2(\C)=4$: one linear condition lowers the dimension by one). We will recognise a nicer basis, $\sx,\sy,\sz$, in Theorem~\ref{thm:paulibasis}.
  \end{steps}
\end{solution}

\begin{pitfall}
  The set of qubit states is \emph{not} a vector space. A qubit state is a \emph{unit} vector of $\C^2$, and the sum of two unit vectors is generally not a unit vector, nor is twice a unit vector. The states live \emph{inside} a vector space ($\C^2$), but they do not form one. This is why normalisation will keep reappearing.
\end{pitfall}

\subsection{Span, independence, bases and dimension}

Fix a vector space $V$. A \term{linear combination} of vectors $v_1,\dots,v_k$ is a sum $\sum_i\lambda_iv_i$ with scalars $\lambda_i$, and the set of all of them is the \term{span} of $v_1,\dots,v_k$. Two properties are of interest.

\begin{definition}[Independence, generators, basis]\label{def:basis}
  Vectors $v_1,\dots,v_k$ are \term{linearly independent} if $\sum_i\lambda_iv_i=0$ holds only when all $\lambda_i=0$. They are a \term{system of generators} of $V$ if their span is all of $V$. They are a \term{basis} if they are both. A space with a finite system of generators is \term{finite-dimensional}.
\end{definition}

Independence says that no vector is redundant; generating says that nothing is missing. A basis is the perfect balance, and it is what lets us use coordinates.

\begin{theorem}[Coordinates]\label{thm:coords}
  If $v_1,\dots,v_n$ is a basis of $V$, every $v\in V$ can be written as $v=\sum_i\lambda_iv_i$ in exactly one way. The numbers $\lambda_i$ are the \term{coordinates} of $v$ in this basis.
\end{theorem}
\begin{proof}
  \begin{steps}
    \item \emph{Existence} is the statement that the $v_i$ generate.
    \item \emph{Uniqueness.} Suppose $\sum_i\lambda_iv_i=\sum_i\mu_iv_i$. Subtracting, $\sum_i(\lambda_i-\mu_i)v_i=0$.
    \item \emph{Independence} forces $\lambda_i=\mu_i$ for every $i$. \qedhere
  \end{steps}
\end{proof}

A fundamental fact, which we use without proof, is that in a finite-dimensional space every basis has the same number of elements. That number is the \term{dimension} $\dim V$. Moreover, any $\dim V$ independent vectors, or any $\dim V$ generating vectors, already form a basis, and a subspace $U\subseteq V$ with $\dim U=\dim V$ equals $V$.
For instance $\dim F^n=n$, $\dim P_n=n+1$ (a basis is $1,x,\dots,x^n$) and the $m\times n$ matrices have dimension $mn$.

From now on, in $\C^2$ we write
\[
  \ket0=\begin{pmatrix}1\\0\end{pmatrix},\quad \ket1=\begin{pmatrix}0\\1\end{pmatrix},\quad \ket{\pm}=\frac{\ket0\pm\ket1}{\sqrt2},\quad \ket{\pm\ii}=\frac{\ket0\pm\ii\ket1}{\sqrt2}.
\]
The symbol $\ket{\,\cdot\,}$ is just a way of writing a vector; its meaning will be explained in Section~\ref{sec:hilbert}. The first basis is the \emph{computational basis}.

\begin{example}\label{ex:qubit-bases}
  Show that $\{\ket+,\ket-\}$ is a basis of $\C^2$, write $\ket0$ in it, and write $\ket{+\ii}$ in it.
\end{example}
\begin{solution}
  \begin{steps}
    \item \emph{Independence.} If $\alpha\ket++\beta\ket-=0$, then in components $(\alpha+\beta)/\sqrt2=0$ and $(\alpha-\beta)/\sqrt2=0$, so $\alpha=\beta=0$.
    \item \emph{Basis.} Two independent vectors in a space of dimension $2$ form a basis.
    \item \emph{Express $\ket0,\ket1$.} Adding the definitions, $\ket++\ket-=\sqrt2\,\ket0$; subtracting, $\ket+-\ket-=\sqrt2\,\ket1$. So $\ket0=\tfrac1{\sqrt2}(\ket++\ket-)$ and $\ket1=\tfrac1{\sqrt2}(\ket+-\ket-)$.
    \item \emph{Express $\ket{+\ii}$.} $\ket{+\ii}=\tfrac1{\sqrt2}\bigl(\ket0+\ii\ket1\bigr)=\tfrac12\bigl((1+\ii)\ket++(1-\ii)\ket-\bigr)$.
    \item \emph{Check.} $\tfrac14\bigl(\abs{1+\ii}^2+\abs{1-\ii}^2\bigr)=1$, as it should be for a unit vector (we will see why in Section~\ref{sec:hilbert}).
  \end{steps}
\end{solution}

\begin{example}[A dependent set of functions]\label{ex:dependent-functions}
  Are $\sin x$, $\cos x$ and $\sin(x+1)$ linearly independent as functions on $\R$?
\end{example}
\begin{solution}
  \begin{steps}
    \item \emph{Expand the third function} with the addition formula (Example~\ref{ex:trig}): $\sin(x+1)=\cos1\,\sin x+\sin1\,\cos x$.
    \item \emph{Rearrange.} $\cos1\cdot\sin x+\sin1\cdot\cos x-1\cdot\sin(x+1)=0$: a vanishing combination with coefficients not all zero.
    \item \emph{Conclude.} The three functions are linearly \emph{dependent}. In fact $\Span\{\sin x,\cos x,\sin(x+1)\}$ is two-dimensional.
  \end{steps}
\end{solution}

\begin{example}[Why quantum computers are hard to simulate]\label{ex:nqubits}
  The state space of $n$ qubits is $(\C^2)^{\otimes n}=\C^{2^n}$ (a tensor product, introduced in Section~\ref{sec:tensor}; for now we only need its dimension, $2^n$). How much memory does a general state of $50$ qubits need, at 16 bytes per complex number?
\end{example}
\begin{solution}
  \begin{steps}
    \item \emph{Count the amplitudes.} A general state is a list of $2^{50}\approx1.1\times10^{15}$ complex numbers.
    \item \emph{Memory.} $2^{50}\cdot16=2^{54}\approx1.8\times10^{16}$ bytes, about $18$ petabytes.
    \item \emph{Scale up.} For $300$ qubits the number $2^{300}\approx2\times10^{90}$ exceeds the number of atoms in the observable universe.
    \item \emph{Moral.} The dimension grows exponentially with the number of qubits, and this is the source both of the power of quantum computing and of the difficulty of simulating it.
  \end{steps}
\end{solution}

\subsection{Sums of subspaces}

Two subspaces $W_1,W_2$ of $V$ can be combined: their \term{sum} $W_1+W_2=\{w_1+w_2\}$ is the smallest subspace containing both. The sum is \term{direct}, written $W_1\oplus W_2$, if every vector of the sum can be written \emph{uniquely} as $w_1+w_2$. This happens exactly when $W_1\cap W_2=\{0\}$. If $V=W_1\oplus W_2$, then $\dim V=\dim W_1+\dim W_2$.

A familiar example: every function $f:\R\to\R$ splits as an even part $\tfrac12(f(x)+f(-x))$ plus an odd part $\tfrac12(f(x)-f(-x))$, and the only function that is both even and odd is $0$, so $\{\text{functions}\}=\{\text{even}\}\oplus\{\text{odd}\}$. We will meet this structure again with projectors, which split a space in the same way.

\subsection{Linear maps}

The maps between vector spaces that matter are those that respect the structure.

\begin{definition}[Linear map]\label{def:linmap}
  Let $V,W$ be vector spaces over $F$. A map $T:V\to W$ is \term{linear} if $T(\lambda v)=\lambda T(v)$ and $T(v+w)=T(v)+T(w)$ for all $\lambda\in F$ and $v,w\in V$. The set of linear maps from $V$ to $W$ is written $L(V,W)$, and $L(V)=L(V,V)$.
\end{definition}

A linear map sends $0$ to $0$, so any map with $T(0)\neq0$ is not linear. The map $T(x,y)=2x+y$ is linear, while $T(x,y)=2x+y-1$ is not, because $T(0,0)=-1$. Differentiation $f\mapsto f'$ and integration $f\mapsto\int_0^1f$ are linear; squaring $f\mapsto f^2$ is not. A linear map is determined by what it does to a basis: if $v=\sum_i\lambda_iv_i$ then $T(v)=\sum_i\lambda_iT(v_i)$.

Linear maps can themselves be added and scaled, so $L(V,W)$ is a vector space, of dimension $\dim V\cdot\dim W$ in finite dimension. Maps can also be composed, and composition of linear maps is linear. It is not commutative in general, and this non-commutativity will be at the heart of quantum mechanics.

\begin{example}[Composition does not commute]\label{ex:noncommuting}
  On $\R^2$ let $T(x,y)=(y,x)$ (swap) and $S(x,y)=(x,0)$ (project to the $x$-axis). Compute $ST$ and $TS$.
\end{example}
\begin{solution}
  \begin{steps}
    \item \emph{$ST$ means: first $T$, then $S$.} $ST(x,y)=S(y,x)=(y,0)$.
    \item \emph{$TS$ means: first $S$, then $T$.} $TS(x,y)=T(x,0)=(0,x)$.
    \item \emph{Compare.} $ST(1,0)=(0,0)$ but $TS(1,0)=(0,1)$: the two compositions differ, and $ST-TS\neq0$.
  \end{steps}
\end{solution}

\begin{lookahead}
  The operator $ST-TS$ is called the \term{commutator} of $S$ and $T$ and written $[S,T]$ (and $\{S,T\}=ST+TS$ is their \term{anticommutator}); operators with $[S,T]=0$ \emph{commute}. The size of the commutator is what the uncertainty principle will measure. Observables are operators; two observables that do not commute (Section~\ref{sec:measurement}) cannot both be sharp in the same state, and the position and momentum operators satisfy the extreme version $\hat x\hat p-\hat p\hat x=\ii\hbar$.
\end{lookahead}

Two subspaces measure what a linear map $T:V\to W$ does.

\begin{definition}[Kernel and range]\label{def:kerran}
  The \term{kernel} (or null space) of $T$ is $\Ker T=\{v\in V:T(v)=0\}$, a subspace of $V$. The \term{range} of $T$ is $\Ran T=\{T(v):v\in V\}$, a subspace of $W$. We say $T$ is \term{injective} if $T(v_1)=T(v_2)$ implies $v_1=v_2$, and \term{surjective} if $\Ran T=W$.
\end{definition}

The kernel is the information that $T$ destroys; the range is what it can reach. The dimension of the range is called the \term{rank} of $T$. The two ideas are linked, as the next two results show.

\begin{theorem}\label{thm:inj-ker}
  A linear map $T$ is injective if and only if $\Ker T=\{0\}$.
\end{theorem}
\begin{proof}
  \begin{steps}
    \item \emph{($\Rightarrow$)} If $T$ is injective and $T(v)=0=T(0)$, then $v=0$.
    \item \emph{($\Leftarrow$)} Suppose $\Ker T=\{0\}$ and $T(v_1)=T(v_2)$. Then $T(v_1-v_2)=0$, so $v_1-v_2=0$. \qedhere
  \end{steps}
\end{proof}

\begin{theorem}[Rank--nullity]\label{thm:ranknull}
  If $V$ is finite-dimensional and $T\in L(V,W)$, then $\dim V=\dim\Ker T+\dim\Ran T$.
\end{theorem}
\begin{proof}
  \begin{steps}
    \item \emph{Build a basis adapted to $T$.} Take a basis $u_1,\dots,u_k$ of $\Ker T$ and extend it to a basis $u_1,\dots,u_k,v_1,\dots,v_r$ of $V$. We claim $T(v_1),\dots,T(v_r)$ is a basis of $\Ran T$; this gives $\dim V=k+r$.
    \item \emph{They generate.} For any $v=\sum a_iu_i+\sum b_jv_j$ we have $T(v)=\sum b_jT(v_j)$, since $T(u_i)=0$.
    \item \emph{They are independent.} If $\sum b_jT(v_j)=0$, then $\sum b_jv_j\in\Ker T$, so $\sum b_jv_j=\sum a_iu_i$ for some $a_i$.
    \item \emph{Conclude.} Independence of the whole basis forces all the $a_i$ and $b_j$ to vanish. \qedhere
  \end{steps}
\end{proof}

\begin{corollary}\label{cor:inj-surj}
  If $\dim V<\infty$ and $T\in L(V)$, then $T$ is injective $\iff$ $T$ is surjective $\iff$ $T$ is invertible.
\end{corollary}
\begin{proof}
  \begin{steps}
    \item By Theorem~\ref{thm:inj-ker}, $T$ is injective iff $\dim\Ker T=0$.
    \item By rank--nullity this holds iff $\dim\Ran T=\dim V$, that is, iff $\Ran T=V$. \qedhere
  \end{steps}
\end{proof}

The next example shows that finite dimension cannot be dropped.

\begin{example}\label{ex:shifts}
  On the space of complex sequences $(z_1,z_2,\dots)$ consider the left shift $\ell(z_1,z_2,z_3,\dots)=(z_2,z_3,\dots)$ and the right shift $\rho(z_1,z_2,z_3,\dots)=(0,z_1,z_2,\dots)$. Are they injective, surjective, bijective? Why does this not contradict Corollary~\ref{cor:inj-surj}?
\end{example}
\begin{solution}
  \begin{steps}
    \item \emph{$\ell$ is not injective.} It forgets the first entry: $\ell(1,0,0,\dots)=0$.
    \item \emph{$\ell$ is surjective.} $(w_1,w_2,\dots)=\ell(0,w_1,w_2,\dots)$.
    \item \emph{$\rho$ is injective.} $\rho(z)=0$ forces all $z_i=0$.
    \item \emph{$\rho$ is not surjective.} Nothing maps to $(1,0,0,\dots)$. Neither shift is bijective.
    \item \emph{No contradiction.} The space is infinite-dimensional, and Corollary~\ref{cor:inj-surj} needs rank--nullity, which needs $\dim V<\infty$.
  \end{steps}
\end{solution}

\begin{example}\label{ex:ranknull}
  Let $T:\R^3\to\R^2$, $T(x,y,z)=(x+y,\,y+z)$. Find $\Ker T$, $\Ran T$ and check rank--nullity.
\end{example}
\begin{solution}
  \begin{steps}
    \item \emph{Kernel.} $T(x,y,z)=0$ means $x=-y$ and $z=-y$, so $\Ker T=\{(-t,t,-t):t\in\R\}=\Span\{(-1,1,-1)\}$, of dimension $1$.
    \item \emph{Range.} $T(1,0,0)=(1,0)$ and $T(0,0,1)=(0,1)$, so $\Ran T=\R^2$, of dimension $2$.
    \item \emph{Check.} $\dim\R^3=3=1+2$.
  \end{steps}
\end{solution}

\begin{figure}[htb]
  \centering
  \begin{tikzpicture}[>=stealth]
    \node[font=\small] at (-4.5,1.15) {$V$};
    \draw[fill=ink!8,draw=ink] (-5.7,-0.6) rectangle (-1.3,1.0);
    \draw[fill=figacc!25,draw=figacc] (-5.5,-0.4) rectangle (-3.6,0.8);
    \draw[fill=navy!20,draw=navy] (-3.6,-0.4) rectangle (-1.5,0.8);
    \node[font=\small,figacc] at (-4.55,0.2) {$\Ker T$};
    \node[font=\small,navy] at (-2.55,0.2) {the rest};
    \draw[->,thick,navy] (-2.55,-0.65) -- (0.9,-0.65) node[midway,below,font=\small] {$T$ is injective here};
    \draw[->,thick,figacc] (-4.55,-0.65)  -- ++(0,-0.3);
    \node[font=\small,figacc] at (-4.55,-1.25) {sent to $0$};
    \node[font=\small] at (3.3,1.15) {$W$};
    \draw[fill=ink!8,draw=ink] (1.3,-0.6) rectangle (5.3,1.0);
    \draw[fill=navy!20,draw=navy] (1.5,-0.4) rectangle (3.4,0.8);
    \node[font=\small,navy] at (2.45,0.2) {$\Ran T$};
    \node[font=\footnotesize] at (4.3,0.2) {unreached};
  \end{tikzpicture}
  \caption{Rank--nullity: the part of $V$ that $T$ kills ($\Ker T$) and the part that it sends one-to-one onto $\Ran T$ add up to all of $V$, so $\dim V=\dim\Ker T+\dim\Ran T$.}
  \label{fig:ranknull}
\end{figure}

\subsection{Matrices}

Linear maps between finite-dimensional spaces become matrices once we pick bases. Let $v_1,\dots,v_n$ be a basis of $V$ and $w_1,\dots,w_m$ a basis of $W$. For each $i$ write $T(v_i)=\sum_jM_{ji}w_j$. The \term{matrix of $T$} in these bases is $M=(M_{ji})$, an $m\times n$ array whose $i$th column holds the coordinates of $T(v_i)$:
\[
  \text{coordinates of }T(v)\ =\ M\cdot\text{coordinates of }v .
\]
This is the one idea to hold on to: \emph{the columns of the matrix are the images of the basis vectors.}

\begin{example}\label{ex:matrices}
  Find the matrix of $D:P_3\to P_3$, $D(p)=p'$, in the basis $1,x,x^2,x^3$, and of the rotation $R_\theta$ of $\R^2$ by an angle $\theta$ in the canonical basis.
\end{example}
\begin{solution}
  \begin{steps}
    \item \emph{Images of the basis of $P_3$.} $D(1)=0$, $D(x)=1$, $D(x^2)=2x$, $D(x^3)=3x^2$.
    \item \emph{Write them as columns of coordinates:}
      \[
        M_D=\begin{pmatrix}0&1&0&0\\0&0&2&0\\0&0&0&3\\0&0&0&0\end{pmatrix}.
      \]
    \item \emph{Check.} Taking four derivatives of a cubic gives $0$, and indeed $M_D^4=0$.
    \item \emph{The rotation.} $R_\theta e_1=(\cos\theta,\sin\theta)$ and $R_\theta e_2=(-\sin\theta,\cos\theta)$, so $M_{R_\theta}=\begin{pmatrix}\cos\theta&-\sin\theta\\ \sin\theta&\cos\theta\end{pmatrix}$.
  \end{steps}
\end{solution}

A matrix depends on the bases chosen, but the map does not. If $P$ is the matrix whose columns are the new basis vectors written in the old basis, the matrix of $T\in L(V)$ changes by $M'=P^{-1}MP$. The next example shows why it is worth choosing the basis well.

\begin{example}\label{ex:changebasis}
  Let $T(x,y)=(x+y,\,x+y)$ on $\R^2$. Find its matrix in the canonical basis and in the basis $b_1=(1,1)$, $b_2=(1,-1)$.
\end{example}
\begin{solution}
  \begin{steps}
    \item \emph{Canonical basis.} $T(e_1)=(1,1)$ and $T(e_2)=(1,1)$, so $M=\begin{pmatrix}1&1\\1&1\end{pmatrix}$.
    \item \emph{New basis.} $T(b_1)=(2,2)=2b_1$ and $T(b_2)=(0,0)=0\,b_2$, so $M'=\begin{pmatrix}2&0\\0&0\end{pmatrix}$.
    \item \emph{Check the formula $M'=P^{-1}MP$} with $P=\begin{pmatrix}1&1\\1&-1\end{pmatrix}$ and $P^{-1}=\tfrac12\begin{pmatrix}1&1\\1&-1\end{pmatrix}$: $MP=\begin{pmatrix}2&0\\2&0\end{pmatrix}$ and $P^{-1}MP=\begin{pmatrix}2&0\\0&0\end{pmatrix}$.
    \item \emph{Reading.} In the good basis the map is simply ``stretch $b_1$ by $2$ and kill $b_2$'' (Figure~\ref{fig:changebasis}).
  \end{steps}
\end{solution}

\begin{figure}[htb]
  \centering
  \begin{tikzpicture}[>=stealth,scale=1.15]
    \draw[ink!20] (-2.4,-2.4) grid (2.4,2.4);
    \draw[->,ink] (-2.4,0)--(2.4,0); \draw[->,ink] (0,-2.4)--(0,2.4);
    \draw[very thick,navy,->] (0,0)--(1,1) node[above left] {$b_1$};
    \draw[very thick,navy,->] (0,0)--(1,-1) node[below right] {$b_2$};
    \draw[very thick,figacc,->] (0,0)--(2,2) node[right] {$T(b_1)=2b_1$};
    \fill[figacc] (0,0) circle (2.2pt);
    \node[figacc,font=\small,anchor=north east] at (-0.1,-0.1) {$T(b_2)=0$};
  \end{tikzpicture}
  \caption{The map $T(x,y)=(x+y,x+y)$ in the basis $b_1,b_2$: it stretches $b_1$ by a factor $2$ and sends $b_2$ to $0$. Its matrix there is $\diag(2,0)$.}
  \label{fig:changebasis}
\end{figure}

\paragraph{Rank, transpose and adjoint.}
The \term{rank} of a matrix is the dimension of the space generated by its columns, which is the same as that generated by its rows, and it equals $\dim\Ran T$. The \term{transpose} $M^{t}$ swaps rows and columns, $(M^t)_{ij}=M_{ji}$. The \term{adjoint} (or conjugate transpose) $M^\dagger$ also conjugates each entry, $(M^\dagger)_{ij}=\overline{M_{ji}}$. For real matrices the two coincide.
All three have the same rank, and $(AB)^t=B^tA^t$, $(AB)^\dagger=B^\dagger A^\dagger$: the order flips.

\paragraph{Square matrices.}
Some square matrices $M$ over $\C$ have names that we will use constantly. The \term{trace} is $\Tr M=\sum_iM_{ii}$, the sum of the diagonal.

\begin{definition}[Special matrices]\label{def:special}
  A square complex matrix $M$ is called
  \begin{itemize}
    \item \term{Hermitian} if $M^\dagger=M$;
    \item \term{unitary} if $M^\dagger M=MM^\dagger=\Id$, that is, $M^{-1}=M^\dagger$;
    \item \term{involutory} if $M^2=\Id$;
    \item \term{idempotent} if $M^2=M$, and a \term{projector} if it is also Hermitian;
    \item \term{nilpotent} if $M^k=0$ for some $k$.
  \end{itemize}
\end{definition}

Hermitian matrices will be the observables of quantum mechanics, and unitary matrices the reversible quantum gates. The determinant, which you already know, tells us when a square matrix can be undone: for $A\in M_n(F)$ the following are all equivalent. $A$ is invertible; $\det A\neq0$; $\rank A=n$; the columns of $A$ are independent; the rows of $A$ are independent; $0$ is not an eigenvalue of $A$ (eigenvalues are the subject of Section~\ref{sec:operators}); the system $Ax=b$ has a unique solution for every $b$.

\begin{example}[Recognising special matrices]\label{ex:special}
  Which of $A=\begin{pmatrix}0&1\\0&0\end{pmatrix}$, $B=\tfrac12\begin{pmatrix}1&1\\1&1\end{pmatrix}$, $C=\begin{pmatrix}1&1\\0&0\end{pmatrix}$ are Hermitian, idempotent, projectors, nilpotent?
\end{example}
\begin{solution}
  \begin{steps}
    \item \emph{$A$.} $A^2=0$: nilpotent. $A^\dagger=\begin{pmatrix}0&0\\1&0\end{pmatrix}\neq A$: not Hermitian, so not a projector.
    \item \emph{$B$.} $B^\dagger=B$ and $B^2=\tfrac14\begin{pmatrix}2&2\\2&2\end{pmatrix}=B$: Hermitian and idempotent, hence a projector (onto $\ket+$).
    \item \emph{$C$.} $C^2=\begin{pmatrix}1&1\\0&0\end{pmatrix}=C$: idempotent. But $C^\dagger\neq C$, so it is not a projector: it projects onto the $x$-axis \emph{obliquely}.
  \end{steps}
\end{solution}

\subsection{The Pauli matrices}

Three $2\times2$ complex matrices are the workhorses of quantum computing:
\[
  \sx=\begin{pmatrix}0&1\\1&0\end{pmatrix},\qquad \sy=\begin{pmatrix}0&-\ii\\\ii&0\end{pmatrix},\qquad \sz=\begin{pmatrix}1&0\\0&-1\end{pmatrix}.
\]
They are called the \term{Pauli matrices}; in a quantum circuit they are the gates $X$, $Y$ and $Z$. We write $\sigma_1,\sigma_2,\sigma_3$ for $\sx,\sy,\sz$.

\begin{proposition}\label{prop:pauli}
  Each Pauli matrix is Hermitian, unitary, involutory and traceless, with determinant $-1$. Moreover, for $j,k\in\{1,2,3\}$,
  \[
    \sigma_j\sigma_k=\delta_{jk}\Id+\ii\sum_{l=1}^{3}\varepsilon_{jkl}\,\sigma_l,
  \]
  where $\varepsilon_{jkl}$ is $1$ for $(j,k,l)$ an even permutation of $(1,2,3)$, $-1$ for an odd one, and $0$ if two indices coincide.
\end{proposition}
\begin{proof}
  \begin{steps}
    \item \emph{Hermitian.} Clear for $\sx,\sz$; for $\sy$, swapping rows with columns gives $\begin{pmatrix}0&\ii\\-\ii&0\end{pmatrix}$ and conjugating returns $\sy$.
    \item \emph{Trace and determinant.} The traces are $0$ and the determinants are $-1$, $-(-\ii)(\ii)=-1$, $-1$.
    \item \emph{Products.} $\sx^2=\Id$ and $\sx\sy=\begin{pmatrix}0&1\\1&0\end{pmatrix}\begin{pmatrix}0&-\ii\\\ii&0\end{pmatrix}=\begin{pmatrix}\ii&0\\0&-\ii\end{pmatrix}=\ii\sz$; the other products are computed in the same way.
    \item \emph{Unitary.} Each $\sigma_j$ is Hermitian with $\sigma_j^2=\Id$, so $\sigma_j^\dagger\sigma_j=\Id$. \qedhere
  \end{steps}
\end{proof}

In words: each Pauli matrix squares to the identity, and distinct ones multiply to $\pm\ii$ times the third, with $\sx\sy=\ii\sz$, $\sy\sz=\ii\sx$, $\sz\sx=\ii\sy$ and the opposite sign in the reverse order. In particular $\sigma_j\sigma_k=-\sigma_k\sigma_j$ for $j\neq k$: Pauli matrices \emph{anticommute}. This gives short proofs of identities that look laborious by hand.

\begin{example}\label{ex:hadamard}
  The \emph{Hadamard gate} is $H=\frac{1}{\sqrt2}\begin{pmatrix}1&1\\1&-1\end{pmatrix}$. Show that $H=(\sx+\sz)/\sqrt2$, that $H^2=\Id$, and that $H\sx H=\sz$.
\end{example}
\begin{solution}
  \begin{steps}
    \item \emph{The identity.} $\sx+\sz=\begin{pmatrix}1&1\\1&-1\end{pmatrix}$, which gives the first claim.
    \item \emph{$H^2$.} Use $\sx^2=\sz^2=\Id$ and $\sx\sz+\sz\sx=0$: $H^2=\tfrac12(\sx+\sz)^2=\tfrac12(\sx^2+\sx\sz+\sz\sx+\sz^2)=\tfrac12(\Id+0+\Id)=\Id$.
    \item \emph{$H\sx H$.} $H\sx H=\tfrac12(\sx+\sz)\sx(\sx+\sz)=\tfrac12(\Id+\sz\sx)(\sx+\sz)$.
    \item \emph{Expand.} $\tfrac12\bigl(\sx+\sz+\sz\sx\sx+\sz\sx\sz\bigr)=\tfrac12\bigl(\sx+\sz+\sz-\sx\bigr)=\sz$, because $\sz\sx\sx=\sz$ and $\sz\sx\sz=-\sx\sz\sz=-\sx$.
  \end{steps}
\end{solution}

The identity $H\sx H=\sz$ says that the Hadamard gate exchanges the roles of $\sx$ and $\sz$. We will see why when we study eigenvectors.

\begin{example}[Products of Pauli matrices]\label{ex:pauli-products}
  Compute $\sx\sy\sz$ and $\sy\sx$, and show that $\sx\sy\sz=\ii\Id$.
\end{example}
\begin{solution}
  \begin{steps}
    \item \emph{First product.} $\sx\sy=\ii\sz$ (Proposition~\ref{prop:pauli}).
    \item \emph{Multiply by $\sz$.} $\sx\sy\sz=\ii\sz\sz=\ii\Id$.
    \item \emph{The reversed product.} $\sy\sx=-\sx\sy=-\ii\sz$ by anticommutation.
  \end{steps}
\end{solution}

Finally, the Pauli matrices together with the identity form a particularly useful basis.

\begin{theorem}[Pauli basis]\label{thm:paulibasis}
  $\{\Id,\sx,\sy,\sz\}$ is a basis of $M_2(\C)$ over $\C$. If we write $\sigma_0=\Id$, then every $A\in M_2(\C)$ satisfies $A=\tfrac12\sum_{\mu=0}^{3}\Tr(\sigma_\mu A)\,\sigma_\mu$.
\end{theorem}
\begin{proof}
  \begin{steps}
    \item \emph{Generate.} Let $A=\begin{pmatrix}a&b\\c&d\end{pmatrix}$. A direct check shows $A=\tfrac{a+d}{2}\Id+\tfrac{b+c}{2}\sx+\tfrac{\ii(b-c)}{2}\sy+\tfrac{a-d}{2}\sz$, so the four matrices generate $M_2(\C)$.
    \item \emph{Basis.} Since $\dim M_2(\C)=4$, four generators form a basis.
    \item \emph{Coefficients.} If $A=\sum_\nu c_\nu\sigma_\nu$, multiplying by $\sigma_\mu$ and taking the trace gives $\Tr(\sigma_\mu A)=\sum_\nu c_\nu\Tr(\sigma_\mu\sigma_\nu)=2c_\mu$, because $\Tr(\sigma_\mu\sigma_\nu)=2\delta_{\mu\nu}$ by Proposition~\ref{prop:pauli}, as $\Tr\Id=2$ and the $\sigma_l$ are traceless. \qedhere
  \end{steps}
\end{proof}

\begin{example}\label{ex:paulidecomp}
  Decompose $A=\begin{pmatrix}1&2\\3&-1\end{pmatrix}$ and $\ket0\!\bra1=\begin{pmatrix}0&1\\0&0\end{pmatrix}$ in the Pauli basis.
\end{example}
\begin{solution}
  \begin{steps}
    \item \emph{$A$: read off $a=1,b=2,c=3,d=-1$.} The coefficients are $\tfrac{a+d}{2}=0$, $\tfrac{b+c}{2}=\tfrac52$, $\tfrac{\ii(b-c)}{2}=-\tfrac\ii2$, $\tfrac{a-d}{2}=1$.
    \item \emph{Assemble.} $A=\tfrac52\sx-\tfrac\ii2\sy+\sz$.
    \item \emph{$\ket0\!\bra1$: $a=d=c=0$, $b=1$.} The coefficients give $\tfrac12\sx+\tfrac\ii2\sy=\tfrac12(\sx+\ii\sy)$.
    \item \emph{Check.} $\sx+\ii\sy=\begin{pmatrix}0&2\\0&0\end{pmatrix}$.
  \end{steps}
\end{solution}

\begin{example}[The real dimension of the Hermitian matrices]\label{ex:hermitian-dim}
  Show that the Hermitian $2\times2$ matrices form a real vector space of dimension $4$, with basis $\Id,\sx,\sy,\sz$.
\end{example}
\begin{solution}
  \begin{steps}
    \item \emph{It is a real vector space} (not a complex one): sums and \emph{real} multiples of Hermitian matrices are Hermitian, but $\ii M$ is not Hermitian when $M\neq0$.
    \item \emph{General form.} Hermitian means $a,d\in\R$ and $c=\bar b$, so $M=\begin{pmatrix}a&b\\\bar b&d\end{pmatrix}$, with $b=\beta_1-\ii\beta_2$.
    \item \emph{Pauli expansion.} By Theorem~\ref{thm:paulibasis}, $M=\tfrac{a+d}2\Id+\beta_1\sx+\beta_2\sy+\tfrac{a-d}2\sz$, with all four coefficients real.
    \item \emph{Dimension.} The four matrices are independent over $\R$, so the real dimension is $4$: a Hermitian qubit observable is described by four real numbers, the coefficient of $\Id$ and a vector $\vec\beta\in\R^3$ of coefficients of $\sx,\sy,\sz$. (This is the reason that $\rho=\tfrac12(\Id+\vec r\cdot\vec\sigma)$ will describe all qubit states, Section~\ref{sec:density}.)
  \end{steps}
\end{solution}

\begin{exerciseset}
  \exercise{1} Decide which of these are vector spaces (over $\R$ or $\C$, as indicated): $\{(x,y)\in\R^2:xy\ge0\}$; the $2\times2$ complex matrices of trace $0$; the Hermitian $2\times2$ matrices over $\C$, and over $\R$; the $2\times2$ unitary matrices; polynomials of degree exactly $3$; $\{f:[0,1]\to\R\text{ differentiable}: f(0)=0\}$ and $\{f:f(0)=1\}$.

  \exercise{2} Is the set of unit vectors of $\C^2$ closed under addition? Under multiplication by complex scalars? Under multiplication by a scalar of modulus $1$? Explain what this says about the structure of the set of qubit states.

  \exercise{1} Decide whether these sets are linearly independent: $\{(1,-1),(1,2),(2,1)\}\subset\R^2$; $\{(1,\ii),(\ii,-1)\}\subset\C^2$; $\{1,x,x^2,(1-x)^2\}\subset P_2$; $\{e^x,e^{2x},e^{3x}\}$ as functions on $\R$.

  \exercise{1} Show that $\{\ket0,\ket1\}$, $\{\ket+,\ket-\}$ and $\{\ket{+\ii},\ket{-\ii}\}$ are bases of $\C^2$. Express $\ket0$ in each of them.

  \exercise{2} Prove: any set containing $0$ is dependent; $\{v_1,\dots,v_k\}$ is dependent if and only if some $v_j$ is a linear combination of the others; a set $A$ is a basis if and only if every vector is a \emph{unique} linear combination of elements of $A$.

  \exercise{2} Show that $\{v_1,v_2\}$ is independent if and only if $\{v_1+v_2,\,v_1-v_2\}$ is independent.

  \exercise{2} Let $W_1,W_2$ be subspaces of $V$. Show that $W_1\cap W_2$ is a subspace, give an example in $\R^2$ where $W_1\cup W_2$ is not, and prove that $W_1\cup W_2$ is a subspace only if one is contained in the other.

  \exercise{2} Show that $V=W_1\oplus W_2$ if and only if $V=W_1+W_2$ and $W_1\cap W_2=\{0\}$. Prove that the even and the odd functions are subspaces of $\{f:\R\to\R\}$ and that the space is their direct sum.

  \exercise{1} Compute the dimension of: $\C^n$ over $\C$; $M_{m\times n}(\C)$; $P_n$; the diagonal $n\times n$ matrices; the upper triangular ones; the Hermitian $n\times n$ matrices over $\R$; the traceless $n\times n$ complex matrices.

  \exercise{2} Show that the solutions of $y''+y=0$ form a space of dimension $2$ with basis $\{\cos x,\sin x\}$ (you may assume the existence and uniqueness theorem for ODEs), and find another basis of complex exponentials.

  \exercise{1} Decide which of these maps are linear, and specify domain and codomain: $T(x,y)=(xy,x)$; $T(z)=\bar z$ on $\C$ (over $\R$? over $\C$?); $f\mapsto f+1$; $f\mapsto f(0)$; $A\mapsto A^t$; $A\mapsto\Tr A$; $A\mapsto\det A$; $A\mapsto BA-AB$ for a fixed $B$.

  \exercise{2} Let $\{v_1,\dots,v_n\}$ be a basis of $V$ and $w_1,\dots,w_n\in W$ arbitrary. Prove there is a unique linear $T$ with $T(v_i)=w_i$. Find the linear $T:\R^2\to\R^3$ with $T(1,1)=(1,0,2)$, $T(1,-1)=(0,1,0)$ and compute $T(3,1)$.

  \exercise{2} Prove that $L(V,W)$ is a vector space, build an explicit basis when $\dim V=n$, $\dim W=m$, and conclude that $\dim L(V,W)=mn$.

  \exercise{2} Show that the composition of linear maps is linear and associative. On $\R^2$ let $R$ be the rotation by $90^\circ$ and $F(x,y)=(x,-y)$ the reflection in the $x$-axis; compute $RF$ and $FR$, and describe each geometrically.

  \exercise{2} Let $D:P_4\to P_3$, $D(p)=p'$ and $M:P_3\to P_4$, $M(p)=xp$. Compute $DM-MD$ on $P_3$. Which famous commutation relation does this resemble?

  \exercise{1} For each map find kernel, range and dimensions, and verify rank--nullity: $T(x,y,z)=(x-y,\,y-z)$; $D:P_n\to P_n$, $D(p)=p'$; $\Tr:M_n(\C)\to\C$; $T(A)=A-A^t$ on $M_n(\R)$.

  \exercise{2} Let $T\in L(V)$ with $T^2=0$. Prove $\Ran T\subseteq\Ker T$ and deduce $\dim\Ran T\le\frac12\dim V$. Give an example with equality in $\C^2$.

  \exercise{3} Let $T\in L(V)$ with $T^2=T$. Prove $V=\Ker T\oplus\Ran T$. Find an operator for which $\Ker T\cap\Ran T\neq\{0\}$.

  \exercise{1} Find the matrices of: $T(x,y)=(x+2y,\,3x-y,\,y)$ in the canonical bases; the reflection across $y=x$; the orthogonal projection onto the $x$-axis; $T(A)=A^t$ on $M_2(\C)$ in the basis $E_{11},E_{12},E_{21},E_{22}$.

  \exercise{2} Let $T(x,y)=(2x+y,\,x+2y)$. Find its matrix $M$ in the canonical basis and $M'$ in the basis $b_1=(1,1)$, $b_2=(1,-1)$, and check that $M'=P^{-1}MP$. Then prove in general that $\Tr$, $\det$ and rank are invariant under $M\mapsto P^{-1}MP$.

  \exercise{2} The Hadamard gate is the linear map with $H\ket0=\ket+$ and $H\ket1=\ket-$. Using its matrix, show $H\sz H=\sx$ and $H\sy H=-\sy$, and deduce $H\sx H=\sz$ again from $H^2=\Id$ without computing.

  \exercise{1} Classify each of these matrices as Hermitian, unitary, idempotent, projector, nilpotent, involutory (several or none may apply): $\begin{pmatrix}1&0\\0&0\end{pmatrix}$, $\begin{pmatrix}0&0\\1&0\end{pmatrix}$, $\tfrac12\begin{pmatrix}1&-1\\-1&1\end{pmatrix}$, $\tfrac1{\sqrt2}\begin{pmatrix}1&1\\-1&1\end{pmatrix}$, $\begin{pmatrix}1&0\\0&\ii\end{pmatrix}$, $\begin{pmatrix}2&\ii\\-\ii&3\end{pmatrix}$.

  \exercise{2} Prove: the diagonal entries of a Hermitian matrix are real; the general $2\times2$ Hermitian matrix is $\begin{pmatrix}a&c-d\ii\\c+d\ii&b\end{pmatrix}$ with $a,b,c,d\in\R$; if $A,B$ are Hermitian then $AB$ is Hermitian iff $AB=BA$, and $\ii[A,B]$ is Hermitian.

  \exercise{2} Let $P^2=P$. Show that $\Id-P$ is idempotent, that the eigenvalues of $P$ lie in $\{0,1\}$, and find a non-Hermitian idempotent $2\times2$ matrix, describing geometrically what it does.

  \exercise{3} Let $N$ be nilpotent with $N^k=0$. Show that its only eigenvalue is $0$, that $\Id-N$ is invertible with inverse $\Id+N+\dots+N^{k-1}$, and that a nilpotent Hermitian matrix is $0$.

  \exercise{2} Show that the eigenvalues of an involutory matrix are $\pm1$ and that, if $M^2=\Id$, then $P_\pm=\tfrac12(\Id\pm M)$ are idempotents with $P_++P_-=\Id$, $P_+P_-=0$ and $M=P_+-P_-$.

  \exercise{2} Show that $U$ is unitary if and only if its columns form an orthonormal basis of $\C^n$; that products and inverses of unitaries are unitary and $\abs{\det U}=1$; and that the phase gate $S=\diag(1,\ii)$ and $T=\diag(1,e^{\ii\pi/4})$ are unitary with $T^2=S$ and $S^2=\sz$.

  \exercise{1} Compute all nine products $\sigma_j\sigma_k$ and check them against Proposition~\ref{prop:pauli}. Then compute $\sz\sy\sx$ and $(\sx\sy)^2$ using only the rules of the proposition.

  \exercise{2} Define $[A,B]=AB-BA$ and $\{A,B\}=AB+BA$. Show that $[\sigma_j,\sigma_k]=2\ii\sum_l\varepsilon_{jkl}\sigma_l$ and $\{\sigma_j,\sigma_k\}=2\delta_{jk}\Id$, and prove the Jacobi identity $[A,[B,C]]+[B,[C,A]]+[C,[A,B]]=0$ for any matrices.

  \exercise{2} For $a,b\in\R^3$ write $a\cdot\sigma=a_x\sx+a_y\sy+a_z\sz$. Prove $(a\cdot\sigma)(b\cdot\sigma)=(a\cdot b)\Id+\ii(a\times b)\cdot\sigma$. Deduce that for a unit vector $n$ one has $(n\cdot\sigma)^2=\Id$, and that its eigenvalues are $\pm1$.

  \exercise{2} Decompose in the Pauli basis: $\begin{pmatrix}1&3\ii\\-\ii&1\end{pmatrix}$, $\ket0\!\bra0$, $H$ and $\begin{pmatrix}a&b\\c&d\end{pmatrix}$.
\end{exerciseset}
